\section{Introduction}
\label{sec:intro}

In 2015 the United States Office of Personnel Management disclosed that attackers had
exfiltrated the records of 21.5 million people, including 5.6 million
fingerprints~\cite{Gootman16}. The passwords in that breach were reset within days. The
fingerprints could not be: a compromised biometric is compromised permanently. This
asymmetry is why ISO/IEC 24745~\cite{ISO24745} requires three properties of any stored
biometric reference --- \emph{irreversibility}, so the reference cannot be inverted to the
biometric; \emph{unlinkability}, so references for the same subject in different databases
cannot be correlated; and \emph{renewability}, so a compromised reference can be replaced
by a fresh one derived from the same trait.

Systems that keep templates encrypted end-to-end address the first requirement directly.
Blind-Touch~\cite{ChoiWooKim24} is the current reference point: it evaluates the matching
step under CKKS~\cite{CKKS17} so that the server never sees a plaintext template. The
guarantee is strong, and the costs are structural rather than incidental. A 117\,MB Galois
rotation key must be distributed to every client; each query ciphertext is 856\,KB; and
because the scheme packs 512 users into one ciphertext, the database is organised around a
slot-packing ceiling rather than around the workload.

Fuzzy private set intersection offers a different trade. Fuzzy-labelled PSI
(FLPSI)~\cite{BuiCong25} lets a client learn the label attached to any server record
within Hamming distance $\tau$ of its query, at communication linear in the database and
with no key material of consequence. But FLPSI matches \emph{bit-strings}, and a
fingerprint is not a bit-string. Bridging that gap --- turning a noisy image into a short
binary code whose Hamming distance tracks biometric similarity --- is the missing piece.

\paragraph{This paper.} We build the bridge and measure what it costs. A Siamese CNN runs
\emph{on the client, in plaintext}, producing a 16-dimensional embedding; a fixed, public,
training-free Super-Bit projection~\cite{JiEtAl12} binarises it to a 128-bit code; FLPSI
matches the code and returns a label. The design removes the homomorphic layer entirely:
server-side templates are 16 bytes, there is no rotation key, and there is no packing
ceiling.

That is the easy half, and on its own it would be an engineering note. The substance of this
paper is what we found on measuring the composition against the error model its cryptographic
half assumes.

FLPSI's parameters are chosen from a model in which impostor codes are uniform, so the Hamming
distance between unrelated records is $\mathrm{Binomial}(128,\tfrac12)$ with standard deviation
$5.66$. Codes from our bridge --- and, we argue, from any LSH bridge over a learned embedding ---
are not uniform: they are 128 projections of a 16-dimensional vector, so they carry roughly 16
degrees of freedom and inherit the embedding's angular distribution, including the fact that some
fingers genuinely look alike. Over $2.88\times10^{6}$ held-out impostor pairs, the mean impostor distance is $64.0$ --- \emph{exactly} what the model
predicts, because the quantiser balances each bit. The standard deviation is $20.22$,
$3.6\times$ the assumed value. A first-moment sanity check on the codes passes cleanly
while the analysis underlying the parameter choice is wrong by three orders of magnitude:
at the published operating point the per-record false-accept rate is
$4.58\times10^{-2}$ rather than the predicted $2.95\times10^{-5}$.

Nor is this a tuning problem. We sweep 6{,}848 parameter combinations and find none that
reaches even a $10^{-2}$ query-level false-accept rate at $N = 5{,}000$ without rejecting
most legitimate users. The reason is a floor: 8 of our $2.88\times10^{6}$ impostor pairs
produce \emph{identical} 128-bit codes, a rate of $2.8\times10^{-6}$ that no choice of
sub-sampling parameters can go below, since a zero-distance pair is accepted by every
parameterisation. This bound is a property of the code, not of the protocol.

\paragraph{Contributions.}
\begin{description}\itemsep2pt
\item[C1.] A design that decouples perceptual feature extraction (client-side, plaintext)
from private matching (FLPSI), removing the homomorphic-encryption layer, its 117\,MB
rotation key and its slot-packing ceiling, and returning the matched identity rather than
a similarity score (\S\ref{sec:design}).

\item[C2.] A public, training-free quantisation bridge whose interface contract with a
fuzzy-PSI matcher is stated \emph{and measured}: bit balance, inter-bit dependence, and
the empirical sub-sample collision probability against its hypergeometric prediction. We
show the bridge costs $+1.55$ percentage points of EER and that the 16-dimensional
embedding, not the code length, is the accuracy ceiling (\S\ref{sec:design},
\S\ref{sec:eval}).

\item[C3.] A parameter-selection analysis for threshold-PSI biometric matching computed
from the \emph{measured} code distribution rather than a uniform idealisation: the
$1{,}550\times$ false-accept gap, the irreducible collision floor, the measured
$(w,\theta,N)$ frontier, and multi-finger fusion as the only mechanism that crosses the
floor (\S\ref{sec:params}).

\item[C4.] A semi-honest security analysis with an explicit ideal functionality, a formal
leakage function, simulators for both corruptions and a hybrid sequence, together with an
assessment of the stored template against ISO/IEC 24745 --- including two attacks we do
not mitigate and the cost of the mitigations we propose (\S\ref{sec:security}).

\item[C5.] A prototype and its measurements: accuracy with identity-level bootstrap
confidence intervals and DET curves across three degradation levels, a bit-budget-matched
comparison against ITQ, communication and latency at the corrected operating point, and a
two-shard scaling experiment (\S\ref{sec:eval}).
\end{description}

\paragraph{What we do not claim.} We do not claim the measured error rates are acceptable
for $1{:}N$ identification at the database sizes we tested; \S\ref{sec:params} shows they
are not. We do not claim at-rest confidentiality against a compromised server. We do not
claim security against malicious adversaries. We report a composition that is attractive
on storage, key material and scaling structure, and that fails on error rates and template
protection, and we quantify both halves.
